\documentclass[11pt]{article}

% ---------- Packages ----------
\usepackage[a4paper, margin=1in]{geometry}
\usepackage{amsmath, amssymb}
\usepackage{graphicx}
\usepackage{float}
\usepackage{hyperref}
\usepackage{enumitem}
\usepackage{xcolor}
\usepackage{caption}
\usepackage{booktabs}
\usepackage{fancyhdr}

% ---------- Header / Footer ----------
\pagestyle{fancy}
\fancyhf{}
\rhead{Intra IIT Tech Meet 1.0}
\lhead{Autonomous Precision Landing on Dynamic Platforms}
\rfoot{\thepage}

% ---------- Placeholder figure command ----------
% Usage: \placeholderfig{width}{height in cm}{Caption text}{label}
\newcommand{\placeholderfig}[4]{%
\begin{figure}[H]
\centering
\fbox{%
\begin{minipage}[c][#2cm][c]{#1\textwidth}
\centering
\color{gray}\itshape #3
\end{minipage}%
}
\caption{#3}
\label{#4}
\end{figure}
}

\title{\textbf{Autonomous Precision Landing on Dynamic Platforms}\\
\large Vision-Guided Cascaded Control for Quadcopter Landing on a Moving Rover}
\author{Track: UAV Dynamics \& Autonomous Systems}
\date{\today}

\begin{document}

\maketitle

\begin{abstract}
This report presents the design of a real-time autonomous flight controller enabling a quadcopter, simulated in PyBullet, to visually track and land on a ground rover following an unpredictable, continuous trajectory. The system combines ArUco-marker-based visual acquisition, ray-casting-based 3D state estimation refined by a Kalman filter, a cascaded (position $\rightarrow$ attitude $\rightarrow$ motor mixing) PD control architecture, and an altitude-adaptive descent profile. A finite state machine governs the transition from target search through synchronized tracking, controlled descent, and emergency motor cutoff upon touchdown, ensuring compliance with the impact-velocity and tilt-angle safety thresholds.
\end{abstract}

\tableofcontents
\newpage

% =====================================================
\section{Introduction}
% =====================================================

Traditional UAV landing systems rely on static GPS waypoints and fixed landing pads. This problem statement instead requires landing on a \textbf{moving} target -- a ground rover following an unpredictable trajectory -- which requires shifting from static waypoint navigation to dynamic, closed-loop, real-time tracking. The controller must integrate:

\begin{itemize}[leftmargin=2em]
    \item \textbf{Computer vision} for target acquisition (ArUco marker detection),
    \item \textbf{State estimation} to infer relative position, velocity, and acceleration,
    \item \textbf{Multi-loop control} to command pitch, roll, yaw, and thrust,
    \item \textbf{A landing state machine} to safely sequence the mission phases.
\end{itemize}

\placeholderfig{0.7}{6}{Cover illustration: quadcopter descending onto a moving rover in the PyBullet arena}{fig:cover}

% =====================================================
\section{System Architecture}
% =====================================================

The controller is structured as a cascaded pipeline, where each stage transforms its input into a more actionable quantity for the next stage:

\placeholderfig{1.0}{9}{Full system block diagram showing the vision pipeline, state estimator, cascaded PD controllers (position and attitude loops), motor mixer, and feedback paths from drone telemetry}{fig:block-diagram}

% =====================================================
\section{State Estimation Stack}
% =====================================================

\subsection{Visual Acquisition}

The downward-facing camera feed is processed using OpenCV's \texttt{cv2.aruco} module to detect the fiducial marker mounted on the rover, yielding its pixel coordinates and orientation relative to the drone's optical center.

\placeholderfig{0.7}{6}{Sample camera frame showing ArUco marker detection with corner points, marker ID, and computed pixel-center overlay}{fig:aruco-detection}

\subsection{Ray Casting for 3D Localization}

The rover's world-frame position is recovered from the 2D pixel measurement via ray casting, using:
\begin{itemize}[leftmargin=2em]
    \item Drone position (from telemetry),
    \item Drone orientation -- roll, pitch, yaw (from telemetry quaternion),
    \item Camera field of view (FOV),
    \item Marker pixel coordinates.
\end{itemize}

Tilt compensation is essential because the marker appearing at the image center does \emph{not} imply it is directly beneath the drone whenever the drone itself is tilted:

\[
\text{Marker at image center} \;\neq\; \text{Marker directly below drone} \quad \text{(when tilted)}
\]

\placeholderfig{0.7}{6}{Geometric diagram of the ray-casting model: camera frame, tilted optical axis, and projection of the pixel ray onto the ground plane to recover rover world coordinates}{fig:raycast-geometry}

\subsection{Kalman Filter}

To smooth noisy vision measurements and predict rover motion during marker occlusion, a Kalman filter is maintained independently for each horizontal axis. The state vectors are:

\[
\mathbf{x} = \begin{bmatrix} x \\ v_x \\ a_x \end{bmatrix}, \qquad
\mathbf{y} = \begin{bmatrix} y \\ v_y \\ a_y \end{bmatrix}
\]

\textbf{Predict step:} propagate the state forward using a constant-acceleration kinematic model to estimate where the rover \emph{should} be, based on physics alone.

\textbf{Update step:} correct the prediction using the new camera-derived measurement of where the rover \emph{appears} to be.

The filter outputs a fused estimate of rover position, velocity, and acceleration, which is passed downstream to the position controller.

\placeholderfig{0.75}{6}{Kalman filter block diagram showing predict/update cycle, with raw noisy vision measurements versus filtered trajectory plotted over time}{fig:kalman-block}

% =====================================================
\section{Position Controller (Outer Loop)}
% =====================================================

The position error and velocity error between rover and drone are computed as:

\[
e_p = p_{\text{rover}} - p_{\text{drone}}, \qquad
e_v = v_{\text{rover}} - v_{\text{drone}}
\]

A proportional-derivative (PD) law produces the desired acceleration:

\[
a_{\text{des}} = K_p \, e_p + K_d \, e_v
\]

The output of this stage is a desired \emph{acceleration}, not a motor command or attitude angle -- that translation happens in the next stage.

\placeholderfig{0.7}{6}{Plot of position error $e_p$ and commanded acceleration $a_{\text{des}}$ over a representative tracking run, illustrating PD response and settling behavior}{fig:position-pd-response}

% =====================================================
\section{Acceleration to Desired Tilt}
% =====================================================

Using Newton's second law, the desired horizontal force is:

\[
F_x = m \, a_x
\]

Since thrust can only be redirected by tilting the vehicle, the small-angle relationship between desired tilt and desired horizontal acceleration is:

\[
\tan(\theta) \approx \frac{a_x}{g}
\]

Hence:

\[
\text{Desired Acceleration} \;\longrightarrow\; \text{Desired Roll/Pitch}
\]

This mapping is applied per-axis, producing $\phi_{\text{des}}$ (roll, from $a_y$) and $\theta_{\text{des}}$ (pitch, from $a_x$).

% =====================================================
\section{Attitude Controller (Inner Loop)}
% =====================================================

Given desired roll and pitch angles $\theta_{\text{des}}$, and actual attitude and angular velocity from telemetry, the attitude error is:

\[
e_\theta = \theta_{\text{des}} - \theta
\]

A second PD law computes the required torque:

\[
\tau = K_p \, e_\theta + K_d \, e_{\dot\theta}
\]

This inner loop runs at a faster effective bandwidth than the outer position loop, consistent with standard cascaded control design practice, ensuring attitude tracking errors are corrected before they propagate into large positional deviations.

\placeholderfig{0.7}{6}{Attitude controller response plot: desired vs.\ actual roll/pitch angle over time, showing convergence and damping characteristics}{fig:attitude-pd-response}

% =====================================================
\section{Motor Mixer}
% =====================================================

The mixer converts the high-level commands -- total thrust, roll torque, and pitch torque (and yaw torque) -- into four individual motor forces $f_1, f_2, f_3, f_4$, respecting the quadcopter's rotor layout and spin directions.

\placeholderfig{0.6}{6}{Quadcopter motor layout diagram (X or + configuration) with labeled rotors $f_1$--$f_4$, spin directions, and torque arm lengths used in the mixing matrix}{fig:motor-mixer}

% =====================================================
\section{Controller Hierarchy Summary}
% =====================================================

\[
\underbrace{\text{Position Error} \rightarrow \text{Position PD} \rightarrow \text{Desired Acceleration}}_{\text{Outer Loop}}
\]
\[
\underbrace{\text{Desired Acceleration} \rightarrow \text{Physics} \rightarrow \text{Desired Roll/Pitch}}_{\text{Acceleration-to-Attitude Mapping}}
\]
\[
\underbrace{\text{Attitude Error} \rightarrow \text{Attitude PD} \rightarrow \text{Desired Torque}}_{\text{Inner Loop}}
\]
\[
\underbrace{\text{Desired Torque} \rightarrow \text{Mixer} \rightarrow \text{Motor Forces}}_{\text{Actuation}}
\]

This nested structure is a classical \textbf{cascaded controller}, chosen because it decouples the slower translational dynamics from the faster rotational dynamics, simplifying both tuning and stability analysis.

% =====================================================
\section{Landing State Machine}
% =====================================================

The mission is governed by a finite state machine with the following phases:

\[
\text{Target Search} \rightarrow \text{Synchronized Tracking} \rightarrow \text{Controlled Descent} \rightarrow \text{Emergency Motor Cutoff}
\]

\placeholderfig{0.8}{7}{Landing state machine diagram showing all states (Target Search, Track, Descend, Blind Descend, Touchdown/Motor Cutoff), transition conditions, and guard logic}{fig:state-machine}

\subsection{Track}
Track the rover in $x, y$ while maintaining constant altitude.

\subsection{Descend}
Track the rover in $x, y$ while progressively reducing altitude; only the $z$-reference changes.

\subsection{Blind Descend}
If the marker is momentarily lost, the Kalman filter's last known position, velocity, and acceleration are used to continue estimating rover motion, and the same controller continues operating on the predicted state rather than stalling.

\subsection{Emergency Motor Cutoff}
Upon detecting physical contact with the landing pad, all motor thrust is instantly cut to prevent the physics engine from flipping or dragging the aircraft.

% =====================================================
\section{Descent Velocity Profiling}
% =====================================================

A constant descent rate (e.g., always $-0.4\ \text{m/s}$) is suboptimal. Instead, an altitude-adaptive profile is used:

\[
v_z = -k \, h
\]

where $h$ is the current altitude above the rover. Representative behavior:

\begin{table}[H]
\centering
\begin{tabular}{@{}cc@{}}
\toprule
Altitude $h$ & Descent Rate $v_z$ \\
\midrule
5\ m    & $-0.5$\ m/s \\
1\ m    & $-0.5$\ m/s \\
0.5\ m  & $-0.25$\ m/s \\
0.1\ m  & $-0.05$\ m/s \\
\bottomrule
\end{tabular}
\caption{Representative altitude-adaptive descent velocity profile}
\label{tab:descent-profile}
\end{table}

\textbf{Rationale:}
\begin{itemize}[leftmargin=2em]
    \item \textbf{Far from the rover:} prioritize descending quickly to close altitude.
    \item \textbf{Near touchdown:} prioritize accurate $x,y$ tracking, since more time is available to eliminate the final landing offset while impact velocity naturally decreases -- directly addressing the $-0.65\ \text{m/s}$ impact-velocity safety threshold.
\end{itemize}

\placeholderfig{0.7}{6}{Plot of descent velocity $v_z$ vs.\ altitude $h$ showing the altitude-adaptive profile compared against a naive constant-rate descent}{fig:descent-profile-plot}

% =====================================================
\section{Edge Case Handling}
% =====================================================

\begin{itemize}[leftmargin=2em]
    \item \textbf{Marker occlusion / loss of visual lock:} handled via Blind Descend using Kalman-filter-predicted rover state.
    \item \textbf{Excessive tilt near touchdown:} attitude PD gains are tuned to stay within the $30^\circ$ chassis tilt limit throughout descent.
    \item \textbf{High impact velocity:} mitigated via the altitude-adaptive descent profile, which drives $v_z \to 0$ as $h \to 0$.
    \item \textbf{Rover sudden maneuvers:} the Kalman filter's acceleration state and the D-term of the position controller jointly compensate for abrupt trajectory changes.
\end{itemize}

% =====================================================
\section{Results}
% =====================================================

\placeholderfig{0.85}{7}{Screenshot / composite of the PyBullet arena showing the full tracking-to-landing trajectory of the quadcopter relative to the moving rover}{fig:results-trajectory}

\placeholderfig{0.7}{6}{Time-series plot of Z-axis velocity at touchdown across multiple test runs, showing compliance with the $-0.65$ m/s impact threshold}{fig:results-impact-velocity}

% =====================================================
\section{Conclusion}
% =====================================================

The presented controller integrates ArUco-based vision, Kalman-filtered state estimation, a cascaded PD control architecture, and an altitude-adaptive descent profile within a landing state machine. In summary:

\begin{quote}
\itshape
Estimate rover position/velocity from camera, compute desired acceleration using PD control, convert acceleration into desired attitude, use another PD controller to achieve that attitude, and mix thrust/torques into motor forces while descending safely onto the rover.
\end{quote}

This layered design isolates each sub-problem -- perception, estimation, translational control, rotational control, and actuation -- enabling independent tuning and robust performance against the dynamic, unpredictable motion of the rover.

\end{document}